% !PRINTED 401
\weberpage{401}

3) Let the given lines be:
\[
0,\ 1,\ 2,\ 3,\ 4,\ 5,\ 5',
\]
and of these let 0, 1, 2, 3, 4 be real, 5 and 5' conjugate imaginary. The
consideration of the three complexes:
\[
\begin{aligned}
(0)\quad & 1\,[1\,0],\ 2\,[2\,0],\ 3\,[3\,0],\ 4\,[4\,0],\ 5\,[5\,0],\ 5'\,[5'\,0] \\
(5)\quad & 0\,[0\,5],\ 1\,[1\,5],\ 2\,[2\,5],\ 3\,[3\,5],\ 4\,[4\,5],\ 5'\,[5'\,5] \\
(5')\quad & 0\,[0\,5'],\ 1\,[1\,5'],\ 2\,[2\,5'],\ 3\,[3\,5'],\ 4\,[4\,5'],\ 5\ [5\,5'],
\end{aligned}
\]
of which the first is real and the two others conjugate imaginary,
and of those derived from (0) by permutation of 0, 1, 2, 3, 4,
gives, exactly as above, the following result:

\begin{center}
\gesperrt{Real double tangents:}
\end{center}
\[
\begin{aligned}
&0,\quad 1,\quad 2,\quad 3,\quad 4,\quad [0\,1],\ [0\,2],\ [0\,3] \\
&[0\,4]\ [1\,2],\ [1\,3]\ [1\,4],\ [2\,3],\ [2\,4],\ [3\,4],\ [5\,5'].
\end{aligned}
\]

\begin{center}
\gesperrt{Conjugate pairs:}
\end{center}
\[
\begin{aligned}
&5,\ [0\,5],\ [1\,5],\ [2\,5],\ [3\,5],\ [4\,5] \\
&5',\ [0\,5'],\ [1\,5'],\ [2\,5'];\ [3\,5'],\ [4\,5'], % [sic] semicolon printed where a comma is expected
\end{aligned}
\]
and we find a syzygetic triple of complexes:
\[
\begin{aligned}
&0,\quad 1,\quad [0\,2][1\,2],\ [0\,3][1\,3],\ [0\,4][1\,4\,],\ [0\,5][1\,5\,],\ [0\,5'][1\,5'] \\
&2,\quad 3,\quad [0\,2][0\,3],\ [1\,2][1\,3],\ [4\,2][4\,3\,],\ [2\,5][3\,5\,],\ [2\,5'][3\,5'] \\
&[0\,1],\ [2\,3],\ [0\,2][1\,3],\ [0\,3][1\,2],\quad 4\ \ [5\,5'],\quad 5\ \ [4\,5'],\quad 5'\ \ [4\,5],
\end{aligned}
\]
of which each complex contains four further real pairs and one
imaginary double pair. But several such triples can be formed,
since 0, 1, 2, 3, 4 may be permuted arbitrarily.
This is the third case of \S.\,104.

4) That, finally, when we suppose all elements of the given
system of seven to be real, all the remaining double tangents
also turn out real, is an immediate consequence of the rational
representation (\S.\,100).

With this it is shown that the four cases of the preceding
section can indeed all occur. Whether the mode of generation
discussed here is the only possible one — in other
words, whether for every real curve of the fourth order a real
\gesperrt{Aronhold} system exists — this question we must leave
undecided.
